\begin{lemma}
Let $\mathcal H$ have finite edge-index set $E$, let $K$ be a finite
color set, and let $I,J:K\to 2^E$. Fix $a\in K$. Suppose $I(a)$ and
$J(a)$ are independent sets in the line graph of $\mathcal H$, and
$I(b)=J(b)$ for every $b\ne a$. For each hypergraph vertex $x$, put
\[
R_I(x)=\{e\in E:x\in e\text{ and }e\notin I(b)\text{ for every }b\in K\},
\]
and define $R_J(x)$ in the same way. Then
$\bigl||R_I(x)|-|R_J(x)|\bigr|\le1$, with cardinalities viewed as real numbers.
\end{lemma}
\begin{proof}
Let $T$ be the set of edges containing $x$ that belong to none of the
sets $I(b)$ for $b\ne a$. The agreement assumption means that the same
set is obtained using $J(b)$ for $b\ne a$. Consequently
$R_I(x)=T\setminus I(a)$ and $R_J(x)=T\setminus J(a)$.
Any two distinct edges of $T$ intersect at $x$, so they are adjacent
by \noderef{LineGraphOfHypergraph}. Independence therefore implies
$|T\cap I(a)|\le1$ and $|T\cap J(a)|\le1$.
The disjoint partitions of $T$ into its intersection with each of these
sets and its complement give
\[
|R_I(x)|=|T|-|T\cap I(a)|,\qquad
|R_J(x)|=|T|-|T\cap J(a)|.
\]
Their difference is $|T\cap J(a)|-|T\cap I(a)|$, which lies in
$[-1,1]$. This proves the asserted absolute-value bound. In the nibble
experiment these are exactly the residual incident edges, since an edge
is colored if at least one color selects it; the tie-breaking rule does
not change whether it is colored.
\end{proof}
